\documentclass[11pt]{article}
\usepackage{amsmath,amssymb,hyperref,geometry}
\geometry{margin=1in}
\title{Recurring warehouse coordination under staggered link recognition}
\date{}
\begin{document}\maketitle

\paragraph{Result.}
Two crews can continue separate work while they disagree about whether their data link is up, and can resume pooled work without resetting who owns a support or cancelling a load already attached to it. The repair is to use delayed topology observations only to drive requests. The physical link interlock checks current joint admission, a durable dock mailbox transfers unique custody, and authentic settlement evidence is still needed before reusing that custody. We give a sharp service time relative to the actual enabling cuts and a uniform envelope over repeated changes. Losing the physical interlock or the separate dock and settlement channels changes the conclusion.

\section{Plant, observations and service}
The civil plant has two controllers $A,B$, one private load-support right per controller and one pooled right $p$. Any accepted job physically reserves its right for exactly $P\ge1$ ticks and then moves to a final rack. Its evidence appears only on acceptance and reaches each controller through the local rack channel after $Q\ge0$ ticks in the implemented subfamily. The link changes between physically connected and disconnected at arbitrary integer ticks, with at most one change per tick. Each generation names the offered buyer if connected; this is a configured work assignment, not an authorization to issue. Each port privately forwards its own authentic link-change record after a delay of at most $D$. Records can arrive out of order; neither actor sees the other's recognition or reads the current true generation. While disconnected no new cross-link data are passed. A request that crossed while connected may be read later from the receiver-side queue, even across multiple changes; it may be duplicated. Request delivery in a stable connected interval takes $L$ ticks in replay. A port knows only its own delayed notice; an ignored or refused request is not a signal of a new topology.

There is one physically shared dock and its link-side interlock, not a replicated global membership oracle. The interlock can validate the actual connected generation for this one joint resource because the data connector actuates that local interlock. It admits at most one pooled job per connected generation for the appointed buyer. It checks the generation of a new deposit as well. It does not publish membership to controllers or grant capacity on other resources. A durable dock mailbox permits the previously deposited custody token to be collected by its named recipient after $h$ ticks independently of link status; it assigns no job. The dock and the rack channels are not affected by the modeled \emph{data-link} split. Consecutive authenticated IDs do not reset. Controllers cannot read actual physical release or actual global jobs through the accounting representation; readiness is determined by their owned right metadata and the learned prefix marks.

The safety law requires exclusive custody and one physical job per support, admission in the current connected generation for a joint job, and no loss or unauthorized cancellation of accepted work. Its old commitments survive subsequent changes irrespective of later epoch recognition. The standing private duty for each continually supplied crew is to start a local job at least every $P+Q$ ticks and complete every private job within $P$ ticks of acceptance. In a connected epoch obeying the stable-window envelope below the appointed buyer is owed a pooled start by the stated bound and a completion $P$ ticks later. An epoch shorter than the envelope may supply pooled service, but an unaccepted offer does not bind or count as completed service. The topology may change again after acceptance before completion. A stable-window promise is not inferred from an individual controller's belief that the link is up.

\section{Installed rule and safety under repeated overlap}
On a locally received change record, retain only a strictly larger generation. The assigned buyer sends one request on that record. Each actor retains at most the newest read request. The nonbuyer deposits the right only if locally holding it, the request matches its own observed generation, and the interlock admits the deposit. Deposit disables the sender immediately. A matching collect enables the recipient once. A stale duplicate token must not create a second owner; a valid outstanding token already deposited can be collected during a newer epoch, but supplies custody only and does not authorize a job. The buyer issues only with local recognition of its assignment, the right in hand, sufficient certified prefixes for that right's recorded prior issues, and interlock acceptance. A stale request or start proposal can fail the interlock, and refusal creates no job. Private rights run the same evidence check but need no link or pooled permission. Attempts for an old generation do not recur at the buyer after refusal.

Safety is inductive. A new deposit removes the only spender and creates one pending custody token; a matching collect consumes that token and creates one spender. A split or rejoin changes neither the token nor a physical reservation. Issuance follows the prior-use prefix test of \href{https://github.com/eternal-november/eos-research-programmes/pull/358}{PR 358}. Under authentic evidence every previous occupant on the spent right has actually accepted at its rack. New physical occupation and settlement preserve exclusivity, and only rack acceptance retires physical occupation. The gate tests current generation independently of stale local beliefs. Therefore all interleavings of admitted phase order, including rejoin before anyone notices a split, preserve the invariant. No synchronization barrier is imposed on private work. Since every issuance earlier than or simultaneous with a private start settles and is evidenced no later than $P+Q$ after it in this symmetric fixed-delay plant, every prefix needed for that private right becomes available within $P+Q$; its greedy next start meets the local service bound. This claim would need a new bound for heterogeneous private dwell or asymmetric evidence transport.

There are two distinct obligations at settlement. The physical support is released on acceptance; the right cannot be reused until its required previous issues are covered by a genuine locally learned contiguous prefix. A received grant atom publishes same-right and within-issuer cumulative dependencies as in PR 358. A new receipt with unknown issue number is disclosed atomically with its grant bundle; a standalone receipt for an already known issue would be harmless. This makes each replica's certified settled set exact for its visible evidence rather than guessed from the new mode. The joined window width is bounded by the number $B=3$ of rights; the underlying conditional criterion is checked in \href{https://github.com/eternal-november/eos-research-programmes/pull/359}{PR 359}. In this extension the replay checks every joined cut against a separate set-of-facts readout.

\section{Stable-window timing and an indistinguishable boundary}
Set the onset of a stable connected generation to tick zero. Let $d_b$ be when the buyer learns it, $d_o$ when the other controller learns it, $\ell$ the buyer-to-owner request read delay, and $r$ the time at which the initial custody arrangement is usable. The dock is either already held by $i$ (then $r=0$), or a valid in-flight token will be collected by $i$ at $r$. An old token is not cancelled or reissued. Let $E$ be the first cut at which the buyer has learned all prefix evidence needed for the previous spend of the pooled right (include zero when already known), and let $a$ be the actual old physical release cut, zero if free. Authentic evidence implies $E\ge a$ for an old pooled occupant. The exact earliest pooled completion for the prescribed installed interface on that cut is
\[
 t_{\rm done}=P+\max\{d_b,E,a, K\},\qquad
 K=\begin{cases}
 r & i=b,\\
 h+\max\{r,d_o,d_b+\ell\}&i\ne b .
 \end{cases}
\]
The lower bound follows, respectively, from local generation recognition, learned settlement and physical release, and the earliest valid custody transfer. For $i\ne b$ a deposit requires custody, the current observed generation and the buyer's matching request, and the recipient requires the subsequent mailbox wait. It is attained by processing ready deliveries before the start phase and issuing at the first enabled tick. The bound applies to rules that obey this request/consent and custody interface, not to proactive delegation before a matching request or extra evidence/authority sources. It is an exact cutwise formula rather than a presumed sum of indivisible delays: all independent waits overlap. The link need remain in the same generation only until the admission cut, inclusive; after that the actual rack service is owed even if it splits again.

For a uniform envelope take $d_b,d_o\le D$, $\ell\le L$, $h,r\le H$ and all old work and necessary authentic prefix evidence due by $P+Q$ after onset. The last condition follows in the equal-delay plant if any required IDs were issued by onset and all rack receipts are delivered independently of the split. Then each generation staying connected through tick
\[
             \max\{P+Q,H+\max(H,D+L)\}
\]
admits the pooled work by then, and finishes within a further $P$ ticks. This is a sufficient uniform envelope, not a claim that all the individual worst cuts are jointly reachable in the same calendar history. In the replay a token deposited at tick 26 survives a split at 28, when the receiver collects it; its stale generation-9 start is refused, and the same token enables generation-11 work at 30 without replacement.

There is an exact stale-view obstruction without a final interlock or equivalent currently valid permission. Just before the tick-1 admission, an owner receives its generation-1 connected notice, has the right and has an empty support. In two evolutions the same physical split and admission attempt occur at tick 1, but the split is ordered just after the admission phase in one and just before it in the other. The next split notice is delayed in both. The owner has the same available inputs in its decision phase. If the connected case owes an immediate start, starting in both violates the new-generation-only admission law in the other, while waiting in both misses that connected deadline. Finer internal bookkeeping cannot distinguish the two without changing evidence. The interlock accepts the first and rejects the second. In the first the load still settles at 5, not at the split. Unbounded recognition delays likewise preclude a finite service guarantee under this installed consent interface: a required buyer's local input can fail to arrive before any proposed deadline. These impossibilities do not exclude designs with different permission or independent evidence mechanisms.

\section{Composition, state and comparison}
A transition block is an ordered partial state update including remaining dock tokens, old obligations, and unread authenticated events on its outgoing boundary. Consecutive blocks compose when their full boundary states agree; ordinary function composition is associative, so rebracketing the same events with the same cuts gives the same result. Documentary view union is separately associative and idempotent for authenticated records from one execution: it does not imply that swapping actual admission with split, or release with acceptance, preserves service. The order witness above changes an actual physical effect with the same local decision. Before comparing at a later cut, simulated policies must be evaluated on their own realized histories rather than merged from separate counterfactual branches.

For $k$ fixed issuer namespaces, each compressed replica maintains $2k$ counters and at most $B$ tail bits up to indexing and fixed right labels, as in PR 358. Each spender retains at most $k$ prior-use marks per owned right, plus its issuer's $k$ cumulative marks, corresponding to $O(kB+k^2)$ global entries. Active custody or mailbox state uses $B$ slots; active physical work is at most $B$. The control layer adds one local generation record, latest request and attempted generation per actor, and one gate generation plus admitted bit. Finite evidence and motion delays retain only recently owed transport items, e.g. at most $2BQ$ individual rack-copy deliveries queued at once for the two actors. At most one topology change can occur per tick, so outstanding notice/request copies are $O(D+L+8)$ for two actors with at most one extra receiver-queue copy read seven ticks later. This bound assumes explicitly bounded age and multiplicity, not unbounded retransmission storage; the simulator's complete trace and raw evidence copies are audit scaffolding, not retained operational memory. Epoch, issuer, custody generation and time identifiers grow logarithmically with execution length when represented in bits, along with authentic proof costs. An unbounded lifetime for old messages invalidates recycling of bounded tags without a separately proved freshness margin. No compact cryptographic implementation is asserted.

\href{https://skoge.folk.ntnu.no/prost/proceedings/ifac2008/data/papers/1774.pdf}{Paoli, Sartini and Lafortune, 2008} model a diagnosing controller switching between nominal and degraded supervisors after a persistent fault; keeping private work safe without waiting for coordination is analogous to the passive/active separation. Their stated single persistent-fault setup and a common diagnosed switch do not establish our recurring live-custody invariant. A factored supervisor on this \emph{same} model, with the local observations and the same interlock, computes the readiness inequality and same proof. A supervisor given true global membership at each actor would have stronger observations, not a fair comparison. The \href{https://github.com/eternal-november/eos-research-programmes/blob/e1348964f73704212d2d8b48d5982a5ed600a2ff/projects/contextual-composition/foundations/native-whole-scope-feedback/Nonblocking_Recurrent_Assembly_Reading_Route.tex}{pinned recurrent assembly} separates available views, ongoing commitments and current-world response. The dock construction preserves those distinctions and instantiates ordered composition on this bounded deterministic state; it does not claim transport by the pinned native law or comparative proof saving.

\section{Executable evidence and disposition}
The full event simulator uses exact copied PR 358 source to enforce ownership, physical release, receipt and replica transitions. It replays 13 alternating link changes with genuinely staggered notices, lost new requests while disconnected, stale delayed requests and notices, valid in-flight custody across a new split, and simultaneous old and new duties across the recurring trace. It checks exact joined views against set union and rebracketing, not merely quiescent equality. A separate lower-bound calculation agrees with 110,592 stable cut cases evaluated by a distinct phase interpreter and 288 comparisons in the full simulator. In the full order control the local input at the decisive attempt is the same, while event ordering changes whether the offered service becomes an accepted duty. These checks support the written invariants and timing proof but do not establish fault tolerance for the dock, rack channels or unbounded evidence lag. The model is a civil warehouse coordination exercise, not physical validation or a general distributed-consensus or irrevocable-availability claim.
\end{document}
